% Source: Weinberg GR, physical PDF pp. 660--664
%         (printed pp. 635--639; blank printed p. 640 is omitted from the scan).
% Coverage: Appendix, Some Useful Numbers.
% Figures/tables/footnotes: numerical and astronomical reference tables.
% Status: source-reviewed and compile-clean.
% Uncertainties: the surprising source value +2600 km/sec for NGC5055 (M63)
%                and all printed question marks are preserved below.

\clearpage
\begingroup
\renewcommand{\thefootnote}{*}
\section*{Appendix: Some Useful Numbers%
  \texorpdfstring{\protect\footnotemark}{*}}
\addcontentsline{toc}{section}{Appendix: Some Useful Numbers}

\subsection*{Numerical Constants}

\begin{center}
\begin{tabular}{@{}ll@{\qquad}ll@{}}
\(\pi\) & \(=3.1415927\) &
\(1''\) & \(=4.8481\times10^{-6}\ \mathrm{radians}\)\\
\(e\) & \(=2.7182818\) &
\(\ln 10\) & \(=2.3025851\)
\end{tabular}
\end{center}

\subsection*{Physical Constants}

\begingroup
\footnotesize
\renewcommand{\arraystretch}{1.16}
\begin{center}
\begin{tabular}{@{}p{0.32\textwidth}p{0.62\textwidth}@{}}
Speed of light &
  \(c=2.9979250(10)\times10^{10}\ \mathrm{cm\,sec^{-1}}\)\\
Gravitational constant &
  \(G=6.6732(31)\times10^{-8}\ \mathrm{dyn\,cm^2\,g^{-2}}\)\\
&
  \(G/c^2=7.425\times10^{-29}\ \mathrm{cm\,g^{-1}}\)\\
Planck's constant &
  \(\hbar=6.582183(22)\times10^{-16}\ \mathrm{eV\,sec}\)\\
&
  \(\hbar=1.0545919(80)\times10^{-27}\ \mathrm{erg\,sec}\)\\
&
  \(h=2\pi\hbar=6.625\times10^{-27}\ \mathrm{erg\,sec}\)\\
Electron volt &
  \(1\ \mathrm{eV}=1.6021917(70)\times10^{-12}\ \mathrm{erg}\)\\
Electronic charge (unrat.) &
  \(e=4.803250(21)\times10^{-10}\ \mathrm{esu}\)
\end{tabular}
\end{center}
\endgroup

\footnotetext{Taken from ``Review of Particle Properties,'' Particle Data
Group, \emph{Rev. Mod. Phys.} \textbf{43}, No.~2, Part II (1971), and
\emph{Astrophysical Quantities}, by C.~W. Allen (Athlone Press, London,
1955). Where figures are given in parentheses, they indicate the one
standard-deviation uncertainty in the last digits of the main numbers.}
\endgroup

\clearpage
\begingroup
\footnotesize
\renewcommand{\arraystretch}{1.16}
\begin{center}
\begin{tabular}{@{}p{0.32\textwidth}p{0.62\textwidth}@{}}
Fine structure constant &
  \(\alpha=e^2/\hbar c=1/137.03602(21)\)\\
Electron mass &
  \(m_e=9.109558(54)\times10^{-28}\ \mathrm{g}\)\\
&
  \(m_ec^2=0.5110041(16)\ \mathrm{MeV}\)\\
Proton mass &
  \(m_p=1.67\times10^{-24}\ \mathrm{g}\)\\
&
  \(m_pc^2=938.2592(52)\ \mathrm{MeV}\)\\
Neutron mass &
  \(m_nc^2=939.5527(52)\ \mathrm{MeV}\)\\
Rydberg &
  \(m_ee^4/2\hbar^2=13.605826(45)\ \mathrm{eV}\)\\
Thomson cross section &
  \(8\pi e^4/3m_e^2c^4=0.6652453(61)\times10^{-24}\ \mathrm{cm^2}\)\\
Weak coupling constant &
  \(g_\beta c/\hbar^3=1.02\times10^{-5}m_p^{-2}\)\\
Boltzmann constant &
  \(k=1.380622(59)\times10^{-16}\ \mathrm{erg\,K^{-1}}\)\\
&
  \(k^{-1}=11604.85(49)\ \mathrm{K/eV}\)\\
Black-body constant &
  \(a_{\mathrm{rad}}=\pi^2k^4/15c^3\hbar^3
    =7.5641\times10^{-15}\ \mathrm{erg\,cm^{-3}\,K^{-4}}\)\\
Typical stellar mass &
  \((\hbar c/G)^{3/2}m_p^{-2}=3.77\times10^{33}\ \mathrm{g}\)
\end{tabular}
\end{center}

\subsection*{General Astronomical Constants}

\begin{center}
\begin{tabular}{@{}p{0.32\textwidth}p{0.62\textwidth}@{}}
Sidereal year (1900) &
  \(1\ \mathrm{year}=3.1558149984\times10^7\ \mathrm{sec}\)\\
Light year &
  \(1\ \mathrm{light\ year}=9.4605\times10^{17}\ \mathrm{cm}\)\\
Mean earth--sun distance &
  \(1\ \mathrm{a.u.}=1.495985(5)\times10^{13}\ \mathrm{cm}\)\\
Parsec &
  \(1\ \mathrm{pc}=3.0856(1)\times10^{18}\ \mathrm{cm}\)\\
&
  \(=3.2615\ \mathrm{light\ year}\)\\
\multicolumn{2}{@{}l@{}}{Hubble time for a Hubble constant of
  \(100\ \mathrm{km\,sec^{-1}\,Mpc^{-1}}\)}\\
&
  \([100\ \mathrm{km\,sec^{-1}\,Mpc^{-1}}]^{-1}
    =9.78\times10^9\ \mathrm{years}\)\\
Solar mass &
  \(M_\odot=1.989(2)\times10^{33}\ \mathrm{g}\)\\
&
  \(M_\odot G/c^2=1.475\ \mathrm{km}\)\\
Solar radius &
  \(R_\odot=6.9598(7)\times10^5\ \mathrm{km}\)\\
Dimensionless solar surface potential &
  \(M_\odot G/R_\odot c^2=2.12\times10^{-6}\)\\
Solar luminosity &
  \(L_\odot=3.90(4)\times10^{33}\ \mathrm{erg\,sec^{-1}}\)\\
Earth mass &
  \(M_\oplus=5.977(4)\times10^{27}\ \mathrm{g}\)\\
&
  \(M_\oplus G/c^2=0.443\ \mathrm{cm}\)
\end{tabular}
\end{center}
\endgroup

\clearpage
\begingroup
\footnotesize
\renewcommand{\arraystretch}{1.16}
\begin{center}
\begin{tabular}{@{}p{0.43\textwidth}p{0.51\textwidth}@{}}
Earth equatorial radius &
  \(R_\oplus=6.37817(4)\times10^3\ \mathrm{km}\)\\
\multicolumn{2}{@{}l@{}}{Dimensionless earth surface potential}\\
&
  \(M_\oplus G/R_\oplus c^2=6.95\times10^{-10}\)\\
\multicolumn{2}{@{}l@{}}{Acceleration due to gravity at earth's surface}\\
&
  \(g=980.665\ \mathrm{cm\,sec^{-2}}\)\\
\multicolumn{2}{@{}l@{}}{Velocity of earth satellite in low orbit}\\
&
  \(v_s=7.9\ \mathrm{km\,sec^{-1}}\)\\
Mean orbital velocity of earth &
  \(v_\oplus=29.78\ \mathrm{km\,sec^{-1}}\)\\
Lunar mass &
  \(M_{\mathrm{Moon}}=7.35\times10^{25}\ \mathrm{g}\)\\
&
  \(M_{\mathrm{Moon}}G/c^2=5.45\times10^{-3}\ \mathrm{cm}\)\\
Lunar radius &
  \(R_{\mathrm{Moon}}=1738\ \mathrm{km}\)\\
\multicolumn{2}{@{}l@{}}{Dimensionless lunar surface potential}\\
&
  \(M_{\mathrm{Moon}}G/R_{\mathrm{Moon}}c^2=3.14\times10^{-11}\)\\
Mean earth--moon distance &
  \(r_{\mathrm{Moon}}=3.84\times10^5\ \mathrm{km}\)\\
\multicolumn{2}{@{}l@{}}{Apparent luminosity of star with apparent
  bolometric magnitude \(m\)}\\
&
  \(l=2.52\times10^{-5}\ \mathrm{erg\,cm^{-2}\,sec^{-1}}
    \times10^{-2m/5}\)\\
\multicolumn{2}{@{}l@{}}{Absolute luminosity of star with absolute
  bolometric magnitude \(M\)}\\
&
  \(L=3.02\times10^{35}\ \mathrm{erg\,sec^{-1}}\times10^{-2M/5}\)
\end{tabular}
\end{center}

\subsection*{Elements of Planetary Orbits}

\begin{center}
\begin{tabular}{@{}lcccc@{}}
\toprule
Planet & Symbol & \(\begin{array}{c}\text{Period}\\T\;(\text{trop.\ year})\end{array}\)
& \(\begin{array}{c}\text{Semilatus rectum}\\L\;(10^6\ \mathrm{km})\end{array}\)
& \(\begin{array}{c}\text{Eccentricity}\\e\;(\text{in 1900})\end{array}\)\\
\midrule
Icarus  & --        & 1.12     & 51.0   & 0.827\\
Mercury & \mercury  & 0.24085  & 55.46  & 0.205615\\
Venus   & \venus    & 0.61521  & 108.20 & 0.006820\\
Earth   & \earth    & 1.00004  & 149.54 & 0.016750\\
Mars    & \mars     & 1.88089  & 225.95 & 0.093312\\
Jupiter & \jupiter  & 11.86223 & 776.5  & 0.048332\\
Saturn  & \saturn   & 29.45772 & 1423   & 0.055890\\
Uranus  & \uranus   & 84.013   & 2863   & 0.0471\\
Neptune & \neptune  & 164.79   & 4498   & 0.0085\\
Pluto   & \pluto    & 248.4    & 5500   & 0.2494\\
\bottomrule
\end{tabular}
\end{center}
\endgroup

\clearpage
\begingroup
\renewcommand{\thefootnote}{*}
\subsection*{Selected Galaxies%
  \texorpdfstring{\protect\footnotemark}{*}}

\footnotesize
\renewcommand{\arraystretch}{1.07}
\begin{center}
\begin{tabular}{@{}lcccc@{}}
\toprule
Local Group & Type & Distance (Mpc) & \(m_{pg}\) &
  \(cz\) (obs.) \((\mathrm{km/sec})\)\\
\midrule
\multicolumn{5}{@{}l}{\emph{Galaxy and Nearest Neighbors}}\\
Galaxy         & Sb or Sc & --    & --          & --\\
LMC            & Ir or SBc& 0.049 & 0.86        & \(+280\)\\
SMC            & Ir       & 0.058 & 2.86        & \(+167\)\\
Ursa Minor     & dE       & 0.077 & ?           & ?\\
Draco          & dE       & 0.08  & ?           & ?\\
Sculptor       & dE       & 0.09  & 10.5        & ?\\
Fornax         & dE       & 0.13  & 9.1         & \(+40\)\\
Leo I          & dE       & 0.23  & 11.27       & ?\\
Leo II         & dE       & 0.23  & 12.85       & ?\\
NGC6822        & Ir       & 0.52  & 9.21        & \(-40\)\\
\addlinespace
\multicolumn{5}{@{}l}{\emph{Other Members of Local Group}}\\
NGC224(M31)    & Sb       & 0.65  & 4.33        & \(-270\)\\
NGC205         & E6p      & 0.65  & 8.89        & \(-240\)\\
NGC221(M32)    & E2       & 0.65  & 9.06        & \(-210\)\\
NGC147         & dE4      & 0.65  & 10.57       & ?\\
NGC185         & E0       & 0.65  & 10.29       & \(-340\)\\
NGC598(M33)    & Sc       & 0.74  & 6.19        & \(-210\)\\
IC1613         & Ir       & 0.74  & 10.00       & \(-240\)\\
Maffei 1       & E        & \(\sim1\) & \(\sim5.8\) (vis.) & ?\\
\addlinespace
\multicolumn{5}{@{}l}{\emph{Miscellaneous Bright Galaxies}}\\
NGC3031(M81)   & Sb       & 2.0   & 7.85        & \(+80\)\\
NGC3034(M82)   & Scp      & 2.0   & 9.20        & \(+400\)\\
NGC5236(M83)   & Sc       & 2.4   & 7.0         & \(+320\)\\
NGC4826(M64)   & ?        & 3.7   & 9.27        & \(+360\)\\
NGC5128(Cen A) & E0p(R)   & \(\sim4.0\) & 7.87 & \(+260\)\\
NGC4736(M94)   & Sbp      & 4.3   & 8.91        & \(+350\)\\
% VERIFY SOURCE: the scan prints +2600 km/sec for NGC5055 (M63).
NGC5055(M63)   & Sb       & 4.3   & 9.26        & \(+2600\)\\
NGC5194(M51)   & Sb       & 4.3   & 9.26        & \(+550\)\\
NGC5457(M101)  & Sc       & 4.3   & 8.20        & \(+400\)\\
\bottomrule
\end{tabular}
\end{center}

\footnotetext{Distances, types, and magnitudes are mostly taken from the
compilation of S.~van den Bergh, \emph{Observer's Handbook of the Royal
Canadian Astronomical Society}, 1971. Under ``Type,'' E denotes
``elliptical,'' with E0, E1, \ldots, increasingly flat. S denotes
``spiral,'' with S0, Sa, Sb, Sc, \ldots, increasingly open; SB denotes
``barred spiral,'' with SB0, SBa, SBb, SBc, \ldots, increasingly open. Ir
denotes ``irregular''; p denotes ``peculiar''; d denotes ``dwarf''; R denotes
``strong radio source.'' For Maffei 1, see H.~Spinrad et al.,
\emph{Ap. J.} \textbf{163}, L25 (1971).}
\endgroup

\clearpage
\begingroup
\small
\renewcommand{\arraystretch}{1.10}
\subsection*{Messier Galaxies in the Virgo Cluster (Visual Magnitudes)}

\begin{center}
\begin{tabular}{@{}lcccc@{}}
\toprule
Galaxy & Type & Distance (Mpc) & \(m_{pg}\) &
  \(cz\) (obs.) \((\mathrm{km/sec})\)\\
\midrule
NGC4472(M49)   & E4      & \(15\pm5\) & 8.9  & --\\
NGC4579(M58)   & SBb     & \(15\pm5\) & 9.9  & --\\
? NGC4621(M59) & E5      & \(15\pm5\) & 10.3 & --\\
NGC4649(M60)   & E2      & \(15\pm5\) & 9.3  & --\\
NGC4303(M61)   & Sc      & \(15\pm5\) & 9.7  & --\\
NGC4374(M84)   & E?      & \(15\pm5\) & 9.8  & --\\
NGC4382(M85)   & S0      & \(15\pm5\) & 9.5  & --\\
NGC4486(M87)   & E0p(R)  & \(15\pm5\) & 9.3  & \(+1220\)\\
NGC4501(M88)   & Sb      & \(15\pm5\) & 9.7  & --\\
NGC4552(M89)   & E0      & \(15\pm5\) & 10.3 & --\\
NGC4569(M90)   & Sb      & \(15\pm5\) & 9.7  & --\\
? NGC4192(M98) & Sb      & \(15\pm5\) & 10.4 & --\\
NGC4254(M99)   & Sc      & \(15\pm5\) & 9.9  & --\\
NGC4321(M100)  & Sc      & \(15\pm5\) & 9.6  & --\\
NGC4594(M104)  & Sb      & \(15\pm5\) & 8.1  & \(+1020\)\\
\bottomrule
\end{tabular}
\end{center}

\subsection*{Selected Clusters of Galaxies}

\begin{center}
\begin{tabular}{@{}lcc@{}}
\toprule
Cluster & Est. No. Galaxies & \(cz\) \((\mathrm{km/sec})\)\\
\midrule
Virgo         & 2500 & 1150\\
Pegasus I     & 100  & 3800\\
Pisces        & 100  & 5000\\
Cancer        & 150  & 4800\\
Perseus       & 500  & 5400\\
Coma          & 1000 & 6700\\
Hercules      & --   & 10300\\
Pegasus II    & --   & 12800\\
Cluster A     & 400  & 15800\\
Ursa Major I  & 300  & 15400\\
Leo           & 300  & 19500\\
Gemini        & 200  & 23300\\
Cor. Bor.     & 400  & 21600\\
Bo\"otes       & 150  & 39400\\
Ursa Major II & 200  & 41000\\
Hydra         & --   & 60600\\
\bottomrule
\end{tabular}
\end{center}
\endgroup
